%
% Manchester Raspberry Jam Workshop Booklet
%
% This template is designed to try and aide consistency between each booklet.
% Insert files as instructed by comments, things like tables of contents are created by other linked files
%

% Is this document PRINT or WEB format?
% (you can use the toggle to cut content for paper, see etoolbox)
\newif\ifprint
\printfalse

% Page Formatting
\ifprint
	\documentclass[a4paper, twocolumn, twoside, 11pt]{article}
	\usepackage[margin=2cm]{geometry}
	\setlength{\columnsep}{1.25cm}
	\setlength{\parskip}{6pt}
\else
	\documentclass[a4paper, onecolumn, oneside, 11pt]{article}
	\usepackage[margin=3cm]{geometry}
	\setlength{\parskip}{9pt}
\fi



% FONT and text format
\usepackage[utf8x]{inputenc}	
\usepackage[UKenglish]{babel}
\ifprint
	\usepackage[usenames, dvipsnames]{color}				%Font Colour
	\usepackage[colorlinks=false]{hyperref}	%URLs
\else
	\usepackage[T1]{fontenc}
	\usepackage[sfdefault]{roboto}
	\usepackage[T1]{fontenc}
	\usepackage[usenames, dvipsnames]{color}				%Font Colour
	\usepackage[colorlinks=true,linkcolor=black,urlcolor=WildStrawberry]{hyperref}
\fi



% Table of Contents Format
\usepackage{tocloft}
\addtocontents{toc}{\cftpagenumbersoff{subsection}}
\setcounter{tocdepth}{2}
\setcounter{secnumdepth}{2}



% Section spacing
\ifprint
	\usepackage[compact]{titlesec}
\fi



% Listings and Asides
\usepackage{McrRaspJam/common/listings}
\usepackage{McrRaspJam/common/asides}



%Clear page for web version
\newcommand{\webclearpage}{
	\ifprint
	\else
		\clearpage
	\fi
}

% misc.
\usepackage{enumitem}									%List spacing changes
\usepackage[toc,page]{appendix}							%Appendix package
\usepackage{graphicx}									% TOC?
\usepackage{etoolbox}									%Boolean used for print/web switching
\usepackage{fancyvrb}									%Centered verbatim

% Enter the document title here
\newcommand{\workshopTitle}{Workshop 17: \textit{Minecraft: Pi Edition}}

% Enter the author of this workshop
\newcommand{\workshopAuthor}{Jack Kelly}


\begin{document}
	% Title Format
\ifprint
	\title{Manchester Raspberry Jam \\ \workshopTitle}
	\author{}
\else
	\title{
		\begin{center}
			\includegraphics[width=30mm]{McrRaspJam/common/logo-512}
		\end{center}
		\vspace{12pt}
		\workshopTitle
	}
	\author{
		Written by \workshopAuthor
	}
\fi

\date{\vspace{-16pt}}
\maketitle


% Online download location
\ifprint
	\begin{mdframed}[rightline=false, leftline=false]
		\scriptsize
		This booklet is available online at \mbox{\href{https://drive.google.com/open?id=0B_1SFjX_5JrmfnhpX0pPRXl6bmJNal8zdUxMeWZOdjJyZVdzU3V6UnBGdlVIMENtbFFkbVk}{bit.ly/McrRaspJam}}
		\normalsize
	\end{mdframed}
\fi
	
	%Place a SINGLE paragraph summary here
	This workshop teaches the basics of the Python API for \textit{Minecraft: Pi Edition}.
	
	%Difficulty
	\textit{Difficulty: Beginner workshop}

	\ifprint
		\renewcommand{\baselinestretch}{0.75}\normalsize
		\tableofcontents
		\renewcommand{\baselinestretch}{1.0}\normalsize
	\else
		\tableofcontents
	\fi
	
	%
	% Input the main CONTENT below, sans title page or contents.
	% Recommend inputting per section, and adding page breaks here.
	%
	% \webclearpage command is provided, will break page for web format only.
	%
	
	\setcounter{section}{-1}
\section{Introduction}

	Included in Raspbian, the operating system we run on our Raspberry Pi's is a copy of the `Pi Edition' of \textit{Minecraft}. 
	
	This version of Minecraft is based on the creative mode from \textit{Minecraft: Pocket Edition}, but in this version we can use the Python programming language to modify the game world.
	
	%Difficulty
	This is an introductory workshop. It's useful to have programmed in Python before,  but all of the programming concepts will be covered from scratch.
		
	\subsection*{How to use these booklets}

	The aim of these booklets is to help you attempt these workshops at home, and to explain concepts in more detail than at the workshop. You don't need to follow these booklets during the workshop, but you can if you'd like 0to.
	
	%Code Listings
		When you need to make changes to your code, they'll be presented in \textit{listings} like the example below. Some lines may be repeated across multiple listings, so check the line numbers to make sure you're not copying something twice.

	\lstinputlisting[style=Python, title=helloworld.py]{McrRaspJam/017_Minecraft/0_introduction/helloworld.py}
	
	
	Occasionally, a concept will be explained in greater detail in \textit{asides}, like the one below. You can read these as you wish, but they're not required to complete the workshop.
	
\begin{aside}[API]
	We use an \textit{Application Programming interface} called `MCPI' to create Python programs for \textit{Minecraft: Pi Edition}.
	
	The API is a set of Python functions which allow us to communicate with the Minecraft game.
\end{aside}

	\subsection*{What you'll need}
		All the software you need for this workshop is pre-installed on recent versions of Raspbian.
	\subsection*{Everything else}
	
		% Aknowledgements
		These booklets were created using \textrm{\LaTeX}, an advanced typesetting system used for several sorts of books, academic reports and letters.
			
		If you'd like to have a look at using LaTeX, We recommend looking at \TeX studio, which is available on most platforms, and also in the 	Raspbian repository.
		
		% License spiel
		To allow modification and redistribution of these booklets, they are distributed under the \hbox{CC BY-SA 4.0} License.
		Latex source documents are available at \url{http://github.com/McrRaspJam/booklet-workshops}
		
		If you get stuck, find errors or have feedback about these booklets, please email me at:
		\href{mailto:jam@jackjkelly.com}{\texttt{jam@jackjkelly.com}}
		\clearpage
		
	\section{Python \& Minecraft Basics}

	Boot your Raspberry Pi to the desktop. The default login is as follows, if your pi doesn't log in automatically.
	
	\begin{tabular}{rl}
		\textbf{username} & pi \\ 
		\textbf{password} & raspberry
	\end{tabular} 
	
	Once at the desktop, Open IDLE from the application menu under \textbf{Programming $\rightarrow$ Python 3 (IDLE)}. The first window that will open is the Python `Shell' window, where our programs will run.
	
	From the shell window press \textbf{File $\rightarrow$ New File} to open a second window. This is where we will write our programs.
	
	\subsection*{Hello, World!}
	
		The first program we write in a new programming language is "Hello, World!". In python, this takes just one line:
		
		\begin{lstlisting}[style=Python, title=helloworld.py]
print("Hello, World!")
		\end{lstlisting}
		
		We can now run our program by clicking \textbf{Run $\rightarrow$ Run Module}, or pressing \textbf{F5}. After saving the file, your program should then run in the shell window.
		
		\begin{lstlisting}[style=Terminal, numbers=none]
=========== RESTART: ===========
Hello, World!
>>> 
		\end{lstlisting}
	
	\pagebreak[1]
	\subsection*{Loading the Minecraft API}
	
		We need to tell Python that this is a Minecraft program. To do this, we start all of our Minecraft programs with these two lines:
	
		\begin{lstlisting}[style=Python, title=hellominecraft.py]
from mcpi.minecraft import Minecraft
mc = Minecraft.create()
		\end{lstlisting}
	
	\pagebreak[1]
	\subsection*{Hello, Minecraft!}
		
		Let's now do the same program in Minecraft, by posting text to the in-game chat.

		If we use \texttt{print()} like we did in the last program, the text would still appear in the shell window. To post to Minecraft, we need to use one of the API functions:	

		\begin{lstlisting}[style=Python, firstnumber=3]

mc.postToChat("Hello, Minecraft!")
		\end{lstlisting}
		
		You can run your program again, but first, make sure Minecraft is running and in-game, otherwise you'll get a an error.
		
	\pagebreak[1]
	\subsection*{Teleportation}
	
		In computing, we keep bits of data in things called \textit{variables}. This could be a number, some text, or something more complex, like the position of a player in a videogame.
		
		In python, we might use variables as follows:
		
		\begin{lstlisting}[style=Terminal, numbers=none]
>>> a = 3
>>> b = 5
>>> a + b
8
		\end{lstlisting}
		
		for our program, we'll get the player position, and place it in a variable called \texttt{pos}.
		
		\begin{lstlisting}[style=Python, firstnumber=5]
		
pos = mc.player.getTilePos()
		\end{lstlisting}
		
		\texttt{pos} is a collection of three numbers, which can be accessed using \texttt{pos.x}, \texttt{pos.y} and \texttt{pos.z}, we can teleport the player by using

		\begin{lstlisting}[style=Python, firstnumber=7]
mc.player.setPos(pos.x, pos.y+100, pos.z)
		\end{lstlisting}
		
		What happens when you run the program? Try using different x, y and z values and see what happens.
		
	\pagebreak[1]	
	\subsection*{Placing Blocks}
		
		to place a block at a set of coordinates, we can do:
		
		\begin{lstlisting}[style=Python, firstnumber=8]

mc.setBlock(pos.x, pos.y, pos.z, 4)
		\end{lstlisting}
		
		this will turn a block into cobblestone. To change the type of block, change the number 4 into one of the block IDs listed in appendix \autoref{sec:blockids}. 
		
		Removing a block is simple, you just need to set the block ID to air! (0)
		
		
		\webclearpage
		
	\section{Nested Loops} \label{sec:loops}

We've just set a block, but what if we want to set many blocks at once?

In computing, when we want to do many things repetitively, we usually use a loop. A python loop might look like this.

\begin{lstlisting}[style=Terminal, numbers=none]
>>> for i in range(1, 100):
	print(i)
1
2
3
...
\end{lstlisting}

if we did a similar thing in a Minecraft program, we could change multiple blocks at once.

\begin{lstlisting}[style=Python, title=minecraftloop.py, breaklines=true]
from mcpi.minecraft import Minecraft
mc = Minecraft.create()

pos = mc.player.getTilePos()

for i in range(0, 10):
	mc.setBlock(pos.x+i, pos.y, pos.z, 4)
\end{lstlisting}

run this program. you should now have a program that draws 10 block in a horizontal line.

A line is a one-dimensional shape. If we wanted to draw a square, we need to create a loop that works in two dimensions.

This is easy, we just need to put one loop inside another:
\begin{lstlisting}[style=Python, title=minecraftloop.py, breaklines=true, firstnumber=6]
for i in range(0, 10):
	for j in range(0, 10):
		mc.setBlock(pos.x+i, pos.y+j, pos.z, 4)
\end{lstlisting}

We've made our loop two dimensional by \textit{nesting} another loop inside it. Each time the loop runs, it creates its own loop.

This time, the program should draw a square, that extends up towards the sky.

A square is a two-dimensional shape. So, how can you change your program once more to create a three-dimensional loop and a cube of blocks?
		\webclearpage
	
	\section{Swords and TNT}

\subsection*{Making TNT Explode}

You may have noticed that whilst you can place TNT using the block ID \texttt{46}, it doesn't explode when you hit it. To make it explode, we need to tell Minecraft to spawn the special, exploding variety of the TNT block.

Each block in minecraft can have a number called a \textit{data value}. This number does something different for different blocks. Usually, it does nothing, sometimes it sets the direction a sign or staircase1 faces. For TNT, it tells the game whether the block can explode or not.

To make TNT explode, we just need to set its data value to 1.

\begin{lstlisting}[style=Python, numbers=none]
mc.setBlock(pos.x, pos.y, pos.z, 46, 1)
\end{lstlisting}

\pagebreak[1]
\subsection*{Sword Events}

The sword in \textit{Minecraft: Pi Edition} has a special feature, it records all the blocks you hit, and the API allows us to do something with each of these block hits.

This Program will make every block you hit with your sword turn into gold!

\begin{lstlisting}[style=Python, title=minecraftloop.py, breaklines=true]
from time import sleep
from mcpi.minecraft import Minecraft
mc = Minecraft.create()

while True:
	blockhits = mc.events.pollBlockHits()
	
	for hit in blockhits:
		pos = hit.pos
		mc.setBlock(pos.x, pos.y, pos.z, 41)
	
	mc.events.clearAll()
	sleep(0.1)	
\end{lstlisting}

Theres a few new things going on here, so let's break it down.

\texttt{blockhits = mc.events.pollBlockHits()} gets the list of recent block hits from the game. This list could be short or long, so we use a \texttt{for} loop to iterate through the list.

For each hit, we take the position of the hit and store it in \texttt{pos}. Then just like before, we can use this pos to set a block.

Finally, we need to tell the game to clear the events list, we only want to change the block that got hit once.

\begin{itemize}[noitemsep]
	\item how can you make the sword turn things into exploding TNT?
	\item how can you make more than one block change each time you make a hit?
	\item is there anything interesting you could do if you \textit{didn't} clear the event list?
\end{itemize}
	\webclearpage
		
	\begin{appendices}
		
		\section{Minecraft Block IDs} \label{sec:blockids}
			
			\ifprint
				\scriptsize
			\else
				\tiny
			\fi
			\begin{tabular}{rl}
				AIR & 0 \\
				STONE & 1 \\
				GRASS & 2 \\
				DIRT & 3 \\
				COBBLESTONE & 4 \\
				WOOD\_PLANKS & 5 \\
				SAPLING & 6 \\
				BEDROCK & 7 \\
				WATER\_FLOWING & 8 \\
				WATER & 8 \\
				WATER\_STATIONARY & 9 \\
				LAVA\_FLOWING & 10 \\
				LAVA & 10 \\
				LAVA\_STATIONARY & 11 \\
				SAND & 12 \\
				GRAVEL & 13 \\
				GOLD\_ORE & 14 \\
				IRON\_ORE & 15 \\
				COAL\_ORE & 16 \\
				WOOD & 17 \\
				LEAVES & 18 \\
				GLASS & 20 \\
				LAPIS\_LAZULI\_ORE & 21 \\
				LAPIS\_LAZULI\_BLOCK & 22 \\
				SANDSTONE & 24 \\
				BED & 26 \\
				COBWEB & 30 \\
				GRASS\_TALL & 31 \\
				WOOL & 35 \\
				FLOWER\_YELLOW & 37 \\
				FLOWER\_CYAN & 38 \\
				MUSHROOM\_BROWN & 39 \\
				MUSHROOM\_RED & 40 \\
				GOLD\_BLOCK & 41 \\
				IRON\_BLOCK & 42 \\
				STONE\_SLAB\_DOUBLE & 43 \\
				STONE\_SLAB & 44 \\
				BRICK\_BLOCK & 45 \\
				TNT & 46 \\
				BOOKSHELF & 47 \\
				MOSS\_STONE & 48 \\
				OBSIDIAN & 49 \\
				TORCH & 50 \\
				FIRE & 51 \\
				STAIRS\_WOOD & 53 \\
				CHEST & 54 \\
				DIAMOND\_ORE & 56 \\
				DIAMOND\_BLOCK & 57 \\
				CRAFTING\_TABLE & 58 \\
				FARMLAND & 60 \\
				FURNACE\_INACTIVE & 61 \\
				FURNACE\_ACTIVE & 62 \\
				DOOR\_WOOD & 64 \\
				LADDER & 65 \\
				STAIRS\_COBBLESTONE & 67 \\
				DOOR\_IRON & 71 \\
				REDSTONE\_ORE & 73 \\
				SNOW & 78 \\
				ICE & 79 \\
				SNOW\_BLOCK & 80 \\
				CACTUS & 81 \\
				CLAY & 82 \\
				SUGAR\_CANE & 83 \\
				FENCE & 85 \\
				GLOWSTONE\_BLOCK & 89 \\
				BEDROCK\_INVISIBLE & 95 \\
				STONE\_BRICK & 98 \\
				GLASS\_PANE & 102 \\
				MELON & 103 \\
				FENCE\_GATE & 107 \\
				GLOWING\_OBSIDIAN & 246 \\
				NETHER\_REACTOR\_CORE & 247
			\end{tabular}
			\normalsize

		\webclearpage
		\section{Other useful functions}

			\begin{lstlisting}[style=Python, numbers=none]
mc.getBlock(x, y, z)
			\end{lstlisting}
			\vspace{-8pt}
			Get a block ID number for the block at position (x,y,z).
			\vspace{8pt}
		
			\begin{lstlisting}[style=Python, numbers=none]
mc.setBlocks(x0, y0, z0, x1, y1, z1, id, [data])
			\end{lstlisting}			
			\vspace{-8pt}
			Place multiple blocks within two coordinates. shortcut for the nested loops from \autoref{sec:loops}.
			\vspace{8pt}
			
			\begin{lstlisting}[style=Python, numbers=none]
mc.getHeight(x, z)
			\end{lstlisting}
			\vspace{-8pt}
			Get the y coordinate (height) of the highest non-air block at that position.
			\vspace{8pt}
			
			\begin{lstlisting}[style=Python, numbers=none]
mc.saveCheckpoint()
			\end{lstlisting}
			\vspace{-8pt}
			Save the world state.
			\vspace{8pt}
			
			\begin{lstlisting}[style=Python, numbers=none]
mc.restoreCheckpoint()
			\end{lstlisting}
			\vspace{-8pt}
			Restore the world state from when it was last saved.
			\vspace{8pt}
			

\iffalse			
		\section{Function quick-reference}
		
		for a full listing see \href{http://www.stuffaboutcode.com/p/minecraft-api-reference.html}{www.stuffaboutcode.com/p/minecraft-api-reference.html}
			
			\begin{lstlisting}[style=Python, numbers=none]
mc.getBlock(x, y, z)
			\end{lstlisting}
		
			\begin{lstlisting}[style=Python, numbers=none]
mc.getBlockWithData(x, y, z)
			\end{lstlisting}
			
			\begin{lstlisting}[style=Python, numbers=none]
mc.setBlock(x, y, z, id, [data])
			\end{lstlisting}
\fi
		
	\end{appendices}
	
	%\section{Java quick start}

	In your preferred editor, create a new file and copy the following code. Note that Java is case sensitive.
	
	\lstinputlisting[style=Java]{McrRaspJam/013_Objects/1_java/HelloWorld.java}
	
	Save this file as \textbf{HelloWorld.java}, then open a terminal to the same folder, and type the commands listed below. You should get the following output.
	
	\begin{lstlisting}[style=Terminal]
$ javac HelloWorld.java
$ java HelloWorld
Hello, World!
	\end{lstlisting}
	
	Now, compare that to the equivalent Python program:
	
	\begin{lstlisting}[style=Python]
print("Hello, World")
	\end{lstlisting}
	
	Coming from most over languages, Java may seem quite complicated. You'll find out that in Java, we have to use a lot of terms that would be optional or excluded in other languages. Java is a language where seemingly nothing is implied.
	
	\begin{aside}[println()]
		println() is the print function for Java. In the Java language we have to state which library or class this function comes from, which in this case is from the \texttt{System} class. \texttt{out} is the standard output stream, text sent to this stream is printed in the terminal.
	\end{aside}
	
	We'll explain what some of the other lines are soon. For now, each program we write will have these same lines, so we don't need to worry about them just yet.
	
	\subsection{Variables}
	
	In Java, variables must be defined with a type. Common types might include:
	
	\begin{itemize}[nosep]
		\item \texttt{int} -- An integer (whole number).
		\item \texttt{double} -- A double-precision floating point number i.e. a number with a decimal point.
		\item \texttt{boolean} -- A boolean value i.e. true or false.
		\item \texttt{String} -- A String i.e text. Strings are a collection of type \texttt{char}.
	\end{itemize}

	An example of variables being defined, then modified is shown below.
	
	\begin{lstlisting}[style=Java, numbers=none]
int a;
a = 5;
int b = 3;
a = a + b;
	\end{lstlisting}
	
	What would the values stored in \texttt{a} and \texttt{b} be after this code was run?	
	
	\subsection{Loops and Conditional Statements}
	
		If you've done any programming before, you'll know that we frequently use loops, or \textbf{iterative statements} to perform repetitive tasks. The \texttt{for} statement is a loop that will operate a set number of times based on what parameters we provide it.
		
		\begin{lstlisting}[style=Java, firstnumber=5]
for(int i=0; i<10; i++)
{
	// Code to repeat
}
else
		\end{lstlisting}
	
		The above \texttt{for} loop would loop it's contents 10 times, starting with \texttt{i} at 0, and stopping when \texttt{i} reaches 10. \texttt{i++} is a shorthand notation for \texttt{i = i + 1}, so each time the loop ends \texttt{i} is increased by 1.
		
		Let's modify our hello world program to print multiple times using this loop. Make the changes in the listing below, then compile (\texttt{javac}) and run (\texttt{java}) your program as before, and you should now see ``Hello, World!'' printed 10 times.
		
		\lstinputlisting[style=Java, breaklines=true]{McrRaspJam/013_Objects/1_java/HelloWorldLoop.java}
		
		A \textbf{conditional statements} like \texttt{if else} are used to selectively run code based on a condition. Let's use an \texttt{if else} statement to change our program a little. We'll change the \texttt{println()} to print out \texttt{i}, the current count number. Then, when we reach a specific number, we'll print out a bit of text after the number.
		
		\lstinputlisting[style=Java, breaklines=true, firstline=5, firstnumber=5, lastline=11]{McrRaspJam/013_Objects/1_java/HelloWorldIf.java}
		
		Now compile and run the program again, you should see something like
		
		\begin{lstlisting}[style=Terminal, numbers=none]
...
4
5
My favourite number!
6
...
		\end{lstlisting}
		
		\textbf{Optional:} A standard basic programming test is FizzBuzz. Use loops and conditions to recreate the game.
		
		\iffalse		
		If you want to try to make something on your own with these statements, a common programming exercise with loops and conditions is FizzBuzz. You'll need the modulus operator \texttt{\%} for your conditions.
		
		\begin{aside}[FizzBuzz]
			FizzBuzz is a number game that is often used as a basic example of a programming task. The objective is to count up from 1 to a certain number, but for each number that is divisible by 3 say `Fizz', for each number divisible by 5 say `Buzz', and for numbers that are divisible by both numbers, say `FizzBuzz'.
			
			The answer for FizzBuzz starts 1, 2, Fizz, 4, Buzz, Fizz, 7, 8, Fizz, Buzz, 11, Fizz, 13, 14, FizzBuzz, 16...
		\end{aside}
		\fi
	
	\subsection{Methods}
	
		\begin{aside}[Sequential Execution]
			A traditional computer program runs one line at a time from top to bottom. In a language like Java, the \textit{main method} is the thing that runs from start to end, and the program ends when the main method is complete. We use method calls to `extend' the main method by jumping to additional code.
		\end{aside}
	
		Java also has \textbf{methods}, often called \textbf{functions} or \textbf{procedures} in other languages, used in much the same way.
		
		Let's write a method that prints out a square number. In our current program, we'll define a new method like the following listing. We also need to make a call to the function from \texttt{main()}.  When you run your program again, it should now print `Hello, World!' 10 times.
		
		\lstinputlisting[style=Java, breaklines=true, firstline=3, firstnumber=3,  lastline=14]{McrRaspJam/013_Objects/1_java/FunctionA.java}
		
		\textbf{Parameters} are passed in the same way as many other languages, within the parentheses with a type deceleration. Functions can also \textbf{return} a value, but the method deceleration must be changed to say what type it returns.
		
		Currently the function is \texttt{void}, meaning it returns nothing. To return a number, we need to change \texttt{void} to \texttt{int}. Let's make the function return the square of the parameter input, then \texttt{main()} will print that number.
		
		\lstinputlisting[style=Java, breaklines=true, firstline=3, firstnumber=3, lastline=15]{McrRaspJam/013_Objects/1_java/FunctionB.java}
	
	
	%	\webclearpage
	
\end{document}